%%%%%%%%%%%%%%%%%%%%%%%%%%%%%%%%%%%%%%%%%%%%%%%%%%%%%%%%%%%%%%%%%%%%%%
% DGX Spark -- NVIDIA DGX Spark field manual
% Author: Prof. Dr.-Ing. Mark Schutera
% Based on Tufte-style layout
\makeindex
%%%%%%%%%%%%%%%%%%%%%%%%%%%%%%%%%%%%%%%%%%%%%%%%%%%%%%%%%%%%%%%%%%%%%%
%
% A field manual for the lab's NVIDIA DGX Spark. Same shape as
% notes_3dprinting: short chapters, one topic each, every chapter also
% ships as its own handout PDF under chapter_pdfs/.
%
% Chapter spine.
%   01 Quickstart                  -- WRITTEN, from laptop to a shell on the Spark
%   02 Settle In                   -- WRITTEN, commented out, folders, venv, detached runs
%   03 Run on the GPU              -- WRITTEN, commented out, cap memory, start small, free it
%
%%%%%%%%%%%%%%%%%%%%%%%%%%%%%%%%%%%%%%%%%%%%%%%%%%%%%%%%%%%%%%%%%%%%%%

\documentclass{../sharedAssets/tufte}

%%
% Book metadata
\title{DGX Spark}
\author[\defaultauthor]{\defaultauthor}
\publisher[\defaultpublisher]{\defaultpublisher}
\renewcommand{\notesdir}{notes_dgxspark}

% Load optional fonts
\loadoptionalfonts

% Additional packages
\usepackage{amsmath}
\usepackage{soul}
\usepackage{textcase}
\usepackage{enumitem}
\usepackage{tabularx}
\usepackage{fancyvrb}
\usepackage{float}
% Break long URLs anywhere so href'd links in the narrow Tufte margins
% wrap instead of running off the side of the page.
\usepackage{xurl}
\Urlmuskip=0mu plus 1mu\relax

% Let chapters open on any page instead of always on a recto. The shared
% class hardcodes \LoadClass[a5paper]{tufte-book} and forwards no options,
% so the flag is set here rather than passed (as in notes_3dprinting).
\makeatletter
\@openrightfalse
\makeatother

\begin{document}

\makestandardfrontmatter
% Three line epigraph: \openepigraph letterspaces a \\ inside its quote
% argument into a stray indent, so its body is spelled out here, one
% \allcaps per line.
\makededication{%
  \begin{fullwidth}
  \sffamily\large
  \begin{doublespace}
  \noindent\allcaps{``Poca favilla gran fiamma seconda.''}\\
  \noindent\allcaps{``A great flame follows a little spark.''}\\
  \noindent\allcaps{Dante Alighieri, Paradiso I, 34}
  \end{doublespace}
  \end{fullwidth}}
\tableofcontents

% Introduction
\makeintroduction{%
  \newthought{This is a manual.} It gets you onto the lab's NVIDIA DGX
  Spark: A desktop machine with one GB10 Grace Blackwell chip and 128~GB
  of memory shared between CPU and GPU, running DGX OS, NVIDIA's Ubuntu.
  It is here for training runs, not as an inference server.
  This is a powerful compute substrate, use it. Then again you will probably never be
  alone on this machine, so be kind to data, jobs, and artifacts on this machine.
  Sandbox whatever you do, so you don't hamper the work of others. 
  On the other hand, expect that someone else could nuke your stuff, so backups are advised.
  If you put stuff on that machine, expect that someone else will be able to see and use it.
  \vspace{1.5em}
  
  \marginnote{\newthought{The official documentation.} NVIDIA's user
  guide and the playbooks at \url{https://build.nvidia.com/spark}.}
  This is currently a quickstart only, so chapters on how to setup a proper project structure and environment
  is highly appreciated. Same holds for a chapter on how to effectively run jobs on a shared GPU. Anything is worth adding, that you would have loved to have known before diving deep.
  The goal is to have a short, concise, and practical manual for the lab's DGX Spark over time.
  If a step is deprecated or you found a better way, contribute it back.
  }


%% Start the main matter
\mainmatter
\authorinrectoheader

% ----------------------------------------------------------------------
% Short chapters, one topic each.
% ----------------------------------------------------------------------
% !TEX root = ../notes_dgxspark.tex
% Chapter 1: Quickstart -- from your laptop to a shell on the Spark
%%%%%%%%%%%%%%%%%%%%%%%%%%%%%%%%%%%%%%%%%%%%%%%%%%%%%%%%%%%%%%%%%%%%%%
%
% Access path: WireGuard tunnel into the lab network, then ssh. Accounts
% and tunnel configs are issued by Mark by mail. mDNS does not cross the
% tunnel, so the chapter uses the IP address from that mail, never the
% spark-xxxx.local name.

\index{quickstart}\index{DGX Spark}
\chapter{Quickstart.}
\printchapterversion{01_quickstart}

\section{Ignition}

\index{account}\newthought{The Spark has one user account per person.}%
\marginnote{\newthought{This setup is far from perfect.} If you want to
build a job scheduler, per-user VPN configurations, or sandboxes that keep
users apart, reach out.}
Request yours by mail. The reply carries your username and
password, and the setup folder as an attached zip with the WireGuard
configuration and the Spark's IP address.
\marginnote{\newthought{Scan the code} and send a mail stating your name, why you need access, which
course, and what you want to run.\par\medskip
\null\hfill\qrcode[height=1.0cm]{mailto:schutera@dhbw-ravensburg.de?subject=DGX\%20Spark\%20Account\&body=Name\%3A\%20\%0ACourse\%3A\%20\%0AProject\%3A\%20}}

\index{sudo}\index{shared infrastructure}\newthought{Your account is an
admin account} on a shared machine. With \texttt{sudo} you can install what
you need, and also break the machine for everyone. With great power comes great
responsibility: Store nothing sensitive on it, leave other users' files and
processes alone, sandbox your stuff and never touch the system config itself. 
Protect yourself from others, and protect others from yourself. 
While you have your own user, two users being logged in at the same time
will produce lag and slowdowns. If you experience this, please log out and try again later.
\marginnote{\newthought{Remember?} Chapter 1 of
Full Stack Handwerkszeug, Development and Production, covers this split,
along with SSH and Linux servers.
\url{https://material.schutera.com/lecturenotes/notes_missingsemester/notes_missingsemester.pdf}}%
That also means the Spark is used most efficiently when you spend as little time as possible to deploy your jobs and leave, 
not for interactive development. If you need to develop interactively, please do so on your own machine.
\marginnote{\newthought{You learn useful stuff here. qed.}}  


\marginnote{\newthought{Download.} \url{https://www.wireguard.com/install/}}
\index{WireGuard}\newthought{WireGuard opens a tunnel into the lab
network.} Only traffic to the lab goes through it; the rest of your
connection stays as it is. Install the WireGuard client, choose
\texttt{Import tunnel(s) from file}, pick \texttt{dgx-spark.conf} from the
setup folder, and press \texttt{Activate}.
The status should now read
\texttt{Active} and \texttt{Latest handshake} is seconds old. No handshake
means your network blocks the tunnel. 

\index{SSH}\newthought{Then log in} with the IP address from the mail:
\begin{verbatim}
ssh <username>@<spark-ip>
\end{verbatim}

\index{nvidia-smi}\newthought{Before you start a job,} run
\texttt{nvidia-smi} to see who is on the GPU.


% % !TEX root = ../notes_dgxspark.tex
% Chapter 2: Settle In -- folders, environment, detached runs
%%%%%%%%%%%%%%%%%%%%%%%%%%%%%%%%%%%%%%%%%%%%%%%%%%%%%%%%%%%%%%%%%%%%%%
%
% Sandbox: everything a student does lives under their home directory.
% Sandbag: their runs are niced, time-limited, and detached.
% The venv + cu130 wheel route is the one NVIDIA's own LLaMA Factory
% playbook uses; the NGC container route is in its Unsloth playbook.

\index{settle in}
\chapter{Settle In}
\printchapterversion{02_settlein}

\section{Your Home Is Your Sandbox}

\index{home directory}\newthought{Everything you do lives under your home
directory.} Nothing goes into \texttt{/opt}, \texttt{/usr/local}, or
another user's folder, and nothing you do in this chapter needs
\texttt{sudo}.
\begin{verbatim}
mkdir -p ~/projects ~/data ~/runs
\end{verbatim}
\marginnote{\newthought{Code, data, and results apart.} Code in
\texttt{projects}, datasets in \texttt{data}, logs and checkpoints in
\texttt{runs}. Deleting a run then never touches your code.}

\index{disk}\newthought{The disk is shared too.} Check what you hold, and
delete what you no longer need. Model downloads pile up in
\texttt{\textasciitilde/.cache}.
\begin{verbatim}
du -sh ~ ~/.cache/*
\end{verbatim}

\section{One Environment per Project}

\index{venv}\index{PyTorch}\newthought{A virtual environment keeps your
packages yours.} Create one inside the project and install PyTorch for
CUDA 13 into it:
\begin{verbatim}
cd ~/projects/myproject
python3 -m venv .venv
source .venv/bin/activate
pip install torch --index-url \
  https://download.pytorch.org/whl/cu130
\end{verbatim}
\marginnote{\newthought{Never \texttt{sudo pip install}.} It writes into
the system Python that every other user runs on.}
\marginnote{\index{ARM}\newthought{The Spark is an ARM machine.} A package
without an \texttt{aarch64} wheel builds from source or fails. Check
before you plan around it.}

\section{Start, Detach, Leave}

\index{nohup}\index{nice}\newthought{Run jobs detached, at low priority,
with a time limit.} The job keeps running after you log out, yields the
CPU to interactive users, and stops itself after eight hours:
\begin{verbatim}
nohup nice -n 10 timeout 8h \
  python train.py > ~/runs/train.log 2>&1 &
\end{verbatim}
\marginnote{\newthought{Check on it later} with
\texttt{tail -f \textasciitilde/runs/train.log}. Stop it with
\texttt{kill} and the PID from \texttt{ps -u \$USER}.}

% % !TEX root = ../notes_dgxspark.tex
% Chapter 3: Run on the GPU -- without taking all of it
%%%%%%%%%%%%%%%%%%%%%%%%%%%%%%%%%%%%%%%%%%%%%%%%%%%%%%%%%%%%%%%%%%%%%%
%
% The Spark has no dedicated GPU memory: CPU and GPU share 128 GB, and so
% do all users. nvidia-smi therefore prints "Memory-Usage: Not Supported"
% (NVIDIA known issue) but still lists per-process GPU memory; total use
% comes from free -h. The default cap of a quarter is Mark's call.

\index{GPU}
\chapter{Run on the GPU}
\printchapterversion{03_gpu}

\section{One Memory for Everyone}

\index{unified memory}\newthought{CPU and GPU share the same 128~GB.} Every
gigabyte your training takes is missing for everyone else, whether it sits
on the CPU side or the GPU side.
\begin{verbatim}
nvidia-smi    # who runs on the GPU, per process
free -h       # how much memory is left overall
\end{verbatim}
\marginnote{\newthought{Memory-Usage: Not Supported} in \texttt{nvidia-smi}
is expected. There is no separate GPU memory to report. The per-process
list below it still works.}

\section{Cap Your Framework}

\index{memory cap}\newthought{Frameworks take more than they need unless
told otherwise.} Cap yours at a quarter of the memory, 32~GB, before the
first tensor touches the GPU:
\begin{verbatim}
torch.cuda.set_per_process_memory_fraction(0.25)  # PyTorch
export XLA_PYTHON_CLIENT_MEM_FRACTION=0.25        # JAX
\end{verbatim}
\marginnote{\newthought{JAX grabs 75\%} of the memory at startup by
default, whether it needs it or not.}
\marginnote{\newthought{Need more than a quarter?} Ask me first.}

\section{Start Small, Then Scale}

\index{batch size}\index{bfloat16}\index{checkpoint}\newthought{Run a few
steps on a slice of the data first.} Watch \mbox{\texttt{free -h}} while it runs,
then raise the batch size until you reach your cap, not the machine's.
Train in \texttt{bfloat16}, reach a large effective batch through gradient
accumulation, and save checkpoints to \texttt{\textasciitilde/runs} so a
stopped job resumes instead of starting over.
\marginnote{\newthought{bfloat16 in PyTorch.} Wrap the training step in
\texttt{torch.autocast} with \texttt{dtype=torch.bfloat16}.}

\newthought{Memory is freed only when your process ends.} An idle Python
session or notebook kernel keeps holding it. Find yours in
\texttt{nvidia-smi} and end it.


% Register the vendor documentation so the reference machinery has entries
% to resolve even while the book has a single chapter.
\nocite{nvidia2026sparkconnect}
\makestandardbackmatter{spark}

\end{document}
